\documentclass[11pt]{scrartcl}

\usepackage[utf8]{inputenc}
\usepackage[ngerman]{babel}
\usepackage[T1]{fontenc}
\usepackage{graphicx}
\usepackage{amsmath}
\usepackage{float}
\usepackage{caption}
\usepackage{siunitx}
\usepackage{xcolor}
\usepackage{geometry}
\geometry{
    a4paper,
    total={170mm,257mm},
    left=20mm,
    top=20mm,
}
\pagestyle{empty}

\title{Gameshow - Großer Preis}
\subtitle{Mittwoch Abendprogramm}
\author{}
\date{}
\begin{document}

\maketitle \thispagestyle{empty}
\section*{Material}
\begin{itemize}
    \item Plakat
    \item Edding(s)
    \item Fragen/Antworten Zettel
    \item Notizzettel und Stifte für Kinder
    \item Stoppuhr
\end{itemize}
\section*{Grundprinzip}
\begin{table}[H]
    \centering
    \begin{tabular}{|c|c|c|c|c|}\hline
        KjG&Geschichte&Geographie&Olympia&Zufall\\\hline
        10&10&10&10&10\\\hline
        20&20&20&20&20\\\hline
        30&30&30&30&30\\\hline
        40&40&40&40&40\\\hline
        50&50&50&50&50\\\hline
    \end{tabular}
    \caption*{Spielplan}
\end{table}
Es gibt mehrere etwa gleich große Gruppen, die Versuchen möglichst viele Punkte zu erspielen. Dabei gibt es 5 Kategorien mit je 5 Schwierigkeitsstufen. Die Gruppe wählt eine Kategorie und eine Schwierigkeitsstufe aus, und erhält die zugehörige Frage oder Aufgabe. Über das Feld verteilt gibt es Joker und Challanges, wobei ein Joker bedeutet, dass es keine Frage oder Aufgabe gibt und eine Challange bedeutet, dass alle Gruppen gegeneinander antreten und der Gewinner erhält die Punkte, auch wenn er nicht an der Reihe ist. Wurde ein Feld gespiel, wird es dem Team, das die Punkte erhalten hat zugeordnet bzw. gestrichen.
\section*{Position in der Rahmenstory} Dies ist nach dem körperlichen Test ein Test des Verstandes, der wissen und Strategie belohnen soll.


\newpage
\section*{Fragen/Antworten}
\subsection*{KjG}
\begin{table}[H]
    \begin{tabular}{|c|c|c|}\hline
        \textbf{Punkte}&\textbf{Frage}&\textbf{Antwort}\\\hline
        10&Was ist das Maskottchen der KjG?&Drache\\\hline
        20&10 Leiter nennen&Gruppenspezifisch\\\hline
        30&Wann war die Neugründung der KjG Rutesheim?&2006\\\hline
        40&\textcolor{red}{\textbf{Challange:}} Wann war unser erste Zeltlager? &2010\\\hline
        50&\textcolor{green}{\textbf{Joker}}&\textcolor{green}{\textbf{Joker}}\\\hline
    \end{tabular}
\end{table}
\subsection*{Geschichte}
\begin{table}[H]
    \begin{tabular}{|c|c|c|}\hline
        \textbf{Punkte}&\textbf{Frage}&\textbf{Antwort}\\\hline
        10&\textcolor{green}{\textbf{Joker}}&\textcolor{green}{\textbf{Joker}}\\\hline
        20&Wie heißen die Gräber der Pharaonen?&Pyramiden\\\hline
        30&Wie lautet der vollständige Name von Caesar&Gaius Iulius Caesar\\\hline
        40&Welches Land wollte Christoph Kolumbus eigentlich erreichen?&Indien\\\hline
        50&Wann war der erste Weltkrieg?&1914-1918\\\hline
    \end{tabular}
\end{table}
\subsection*{Geographie}
\begin{table}[H]
    \begin{tabular}{|c|c|c|}\hline
        \textbf{Punkte}&\textbf{Frage}&\textbf{Antwort}\\\hline
        10&\textcolor{red}{\textbf{Challange:}} Deutsche Städte nennen (abwechselnd)&Grupenspezifisch\\\hline
        20&5 Hauptstädte der Welt nennen&Gruppenspezifisch\\\hline
        30&\textcolor{green}{\textbf{Joker}}&\textcolor{green}{\textbf{Joker}}\\\hline
        40&Nachbarländer von Spanien &PT, FR, AND, MA, GBZ (Gibraltar)\\\hline
        50&\textcolor{red}{\textbf{Challange:}} Meiste Länder in $\qty{60}{\second}$ aufschreiben&Gruppenspezifisch\\\hline
    \end{tabular}
\end{table}
\subsection*{Olympia}
\begin{table}[H]
    \begin{tabular}{|c|c|c|}\hline
        \textbf{Punkte}&\textbf{Frage}&\textbf{Antwort}\\\hline
        10&Wo fanden die Ol.Sp. 2024 statt?&Paris\\\hline
        20&Wann waren die ersten Olympische Spiele der Neuzeit?&1896\\\hline
        30&\textcolor{red}{\textbf{Challange:}} Wie viel Sportler haben 2024 bei den Ol.Sp. teilgenommen?&\num{11.069}\\\hline
        40& &Laufen, Springen,\\ &Aus welchen Sportarten bestand der Pentathlon? (Antike Ol.Sp.)&Diskus, Speer,\\ & &Ringen\\\hline
        50&Wie viele Länder haben an den ersten Ol.Sp. der N.Z. teilgenommen?&14\\\hline
    \end{tabular}
\end{table}
\subsection*{Zufall}
\begin{table}[H]
    \begin{tabular}{|c|c|c|}\hline
        \textbf{Punkte}&\textbf{Frage}&\textbf{Antwort}\\\hline
        10&Kopfrechnen: $125/(27-22)$&$25$\\\hline
        20&Wie heißen die Kinder von Schafen?&Lamm/Lämmer\\\hline
        30&\textcolor{green}{\textbf{Joker}}&\textcolor{green}{\textbf{Joker}}\\\hline
        40&Wie alt sind die Teilis im Durchschnitt?&\textcolor{violet}{\textbf{t.b.d.}}\\\hline
        50&Welche Farbe hat flüssiger/fester Sauerstoff?&blau\\\hline
    \end{tabular}
\end{table}

\end{document}